problem:  
Let \( (M^{2n}, g, J, \omega) \) be a simply connected complex space form of constant holomorphic sectional curvature \(c\) (\(M=\mathbb{C}P^n_c, \mathbb{C}H^n_c,\) or \(\mathbb{C}^n\)).  
Let \(L^n \subset M\) be a compact, oriented, **Hamiltonian stationary Lagrangian submanifold**, so that its Maslov form  
\[
\alpha_H = \iota_H \omega
\]
is harmonic. Denote by 
\[
E(L) = \int_L |\alpha_H|^2 \, d\mathrm{vol}_L
\]
the *Maslov energy*, and by \(\|\mathrm{II}\|\) the norm of the second fundamental form.

**Conjectural Problem (Analytic-Topological Small-Maslov-Class Problem):**  
Fix dimension \(n\) and curvature \(c>0\). Does there exist a function  
\[
\Psi_n:(0,\infty)\to(0,\infty)
\]
with \(\Psi_n(R)\to0\) as \(R\to\infty\) such that the following holds?

> Suppose \(L^n\subset \mathbb{C}P^n_c\) is compact, Hamiltonian stationary, satisfying  
> (1) \(\sup_L \|\mathrm{II}\| \le R\);   
> (2) \(E(L) \le \Psi_n(R)\,\mathrm{Vol}(L)\).  
> Then \(L\) has *exact* Maslov class \( [\alpha_H] = 0 \).  
> In particular, \(L\) is Hamiltonian isotopic to a minimal Lagrangian submanifold (possibly, the totally geodesic real form \( \mathbb{R}P^n \subset \mathbb{C}P^n_c \)).

Equivalently: does sufficiently small \(L^2\)-Maslov energy, under bounded extrinsic curvature, force a Hamiltonian stationary Lagrangian to be Hamiltonian isotopic to a minimal one?

---

Why is it a "good" problem:  

1. **Nonvacuous curvature-energy interplay:**  
   By introducing the curvature control \(\|\mathrm{II}\|\le R\), the problem avoids the degeneracy that harmonic 1‑forms can be scaled to have arbitrarily small energy. This ensures that smallness of \(E(L)\) carries real geometric meaning, producing a nontrivial analytic–topological question.

2. **Bridging geometric analysis and symplectic topology:**  
   The problem directly links the analytic smallness of a harmonic form (\(\alpha_H\)) to a topological property (exactness, i.e. trivial Maslov class). It demands understanding of how energy estimates derived from Bochner formulas on \(L\) interact with the ambient curvature of \(M\), thus connecting elliptic analysis and symplectic geometry.

3. **Quantitative rigidity towards minimal Lagrangians:**  
   If such a function \(\Psi_n\) exists, it would provide a *quantitative stability criterion* for minimal Lagrangian submanifolds within the Hamiltonian stationary class—analogous to energy gap results for minimal surfaces or harmonic maps.

4. **Conceptually simple yet deeply unexplored:**  
   The statement uses only standard geometric quantities (second fundamental form, L²‑norm of the Maslov form) but proposes a very subtle conclusion involving global topology and variational rigidity. Its analytic simplicity conceals profound interplay among extrinsic curvature control, harmonic forms, and symplectic potential theory.

5. **Potential to generate new tools:**  
   Attempting such an energy–topology implication would likely require new curvature estimates, refined Bochner identities for the mean curvature 1‑form, or calibrated geometry techniques—offering conceptual unification between Hamiltonian stationary theory and global analysis on submanifolds.

Hence, this amended version defines a logically sound, nonvacuous, and conceptually rich geometric-analytic conjecture exemplifying the hallmarks of a “good” research problem.